\clearpage\SourceLeaf{448}\ArithmeticHead{441}{R}{BEMERKUNGEN ZUR REIHENLEHRE. ABSCHNITT VII.}
\begingroup\small
\makeatletter\tagsleft@true\let\veqno\@@leqno\makeatother
\noindent Da nun
\[
\lim_{n\to\infty}\Bigl(1+\frac12\cdot\frac{2n+1}{2n+2}\gamma^2+\frac12\cdot\frac34\cdot\frac{2n+1}{2n+2}\cdot\frac{2n+3}{2n+4}\gamma^4+\cdotp\cdotp\cdotp\Bigr)=(1-\gamma^2)^{-\frac12}
\]
\noindent ist, so ergibt sich dieser Koeffizient für große Werte von \(n\) gleich:
\[
2\gamma^n\frac12\cdot\frac34\cdotp\cdotp\cdotp\frac{2n-1}{2n}=2\gamma^n\frac{(2n)!}{2^{2n}(n!)^2},
\]
\noindent und hieraus folgt mit Benutzung der \textsc{Stirling}schen Näherungsformel
\[n!=n^n e^{-n}\sqrt{2\pi n}\]
\noindent der von \textsc{Gauss} in der Gleichung [27], S. 423, angegebene Näherungswert \(2\gamma^n\frac1{\sqrt{\pi n}}\). Nunmehr findet man mit Benutzung der Gleichung [23] für den Koeffizienten \(C_n\) der Entwicklung (III) S. 439 den in der Gl. [28], oben S. 423, angegebenen Wert, also für den Koeffizienten von \(\sin nM\) in der Entwicklung (IV) der Mittelpunktsgleichung den angenäherten Wert
\begin{equation}\tag*{[28]$'$}
\frac{C_n}{n}=\frac1n\Bigl(\operatorname{tg}\frac\varphi2 e^{\cos\varphi}\Bigr)^n\Bigl(1+\frac43\frac1{\sqrt{2n\pi\cos^3\varphi}}\Bigr).
\end{equation}

\indent In der Tabelle am Schluß des art. [3.], oben S. 424, hat \textsc{Gauss} für einen bestimmten Wert der Exzentrizität den Fehler berechnet, den man begeht, wenn man sich auf die durch die Gleichung [23], oben S. 423, dargestellte Annäherung beschränkt. Die 25 Zahlengruppen beziehen sich auf die 25 allemal um \(7^0\,30'\) wachsenden Werte der mittleren Anomalie \(M\) von \(0^0\) bis \(180^0\); die in einer Gruppe untereinander stehenden Zahlen bedeuten der Reihe nach:
\[
\begin{array}{r@{\,}l}
&\text{den Wert von }M\text{ (fett gedruckt)}\\[3pt]
&\log_{10}\cos\varphi\\[3pt]
&\log_{10}(1-f\cos E)^2\\[3pt]
&\log_{10}\frac12\cos\theta\\[3pt]
\pm&\log_{10}(1-g\cos M)^{-1},\ \text{jenachdem }M\lessgtr90^0,\\[5pt]
&\log_{10}\sqrt{\frac{2\cos\theta}{9\cos^3\varphi}}\\[5pt]
\pm&\log_{10}(1-g\cos M)^{-\frac12},\ \text{jenachdem }M\lessgtr90^0,\\[5pt]
&\dfrac{\cos\varphi}{(1-f\cos E)^2}\\[9pt]
&\dfrac12\dfrac{\cos\theta}{1-g\cos M}\\[9pt]
&\sqrt{\dfrac{2\cos\theta}{9\cos^3\varphi}}\ \dfrac1{\sqrt{1-g\cos M}}\\[9pt]
&\dfrac{\cos\varphi}{r^2}-\dfrac12\dfrac{\cos\theta}{1-g\cos M}-\sqrt{\dfrac{2\cos\theta}{9\cos^3\varphi}}\ \dfrac1{\sqrt{1-g\cos M}},
\end{array}
\]
\noindent wo \(r=1-f\cos E\) ist, so daß also die letzte Zahl den begangenen Fehler angibt. Der dieser Tabelle zugrunde liegende Wert der Exzentrizität ist hiernach
\[f=\sin\varphi=0{,}6550518\]
\noindent entsprechend dem Werte
\begin{equation}\tag*{($\ast$)}\log_{10}f=9{,}8162757-10.\end{equation}
\par\smallskip\noindent X1.\hfill 56
\clearpage\SourceLeaf{449}\ArithmeticHead{442}{L}{ANALYSIS. NACHLASS.}
\noindent Für diesen Wert der Exzentrizität ist auch der am Schluß von art. [3.], oben S. 425, angegebene Zahlenwert \(0{,}5508570\) des Ausdrucks [29] berechnet\textsuperscript{*}). Da der konstante Teil von
\[\frac{\cos\theta}{2(1-g\cos M)}\]
\noindent gleich \(\tfrac12\) ist, so ergibt also die Näherungsformel [23] für den konstanten Teil von \(\dfrac{dv}{dM}\) den Wert \(1{,}0508570\) statt 1. — Auch der Wert \(9{,}9046115\) für den Logarithmus des Ausdrucks [35], oben S. 426, und die Zahlenbeispiele des art. [5.], oben S. 427, sind für eben denselben Wert der Exzentrizität berechnet. Für die Beispiele des art. [5.] zeigt dies schon der dort angegebene Wert
\[\log_{10}kf=9{,}4540600,\]
\noindent den man in der Tat erhält, wenn man dem in (\(\ast\)) gegebenen Werte von \(\log_{10}f\) den Wert
\[\log_{10}k=\log_{10}\log_{10}e=9{,}6377843\]
\noindent hinzuzählt. In bezug auf den art. [5.] wäre noch zu bemerken, daß \(\theta\) hier einfach einen durch die Gleichung [36] erklärten Hilfswinkel bedeutet, also mit dem in den artt. [2.]—[4.] durch \(\theta\) bezeichneten Winkel nichts zu tun hat. Die Ausdrücke [37], [38] ergeben sich, wenn man in \(M=E-f\sin E\) einsetzt
\[E=p+i\log\operatorname{tang}(45^0+\tfrac12\theta)\]
\noindent und dann ausdrückt, daß \(M\) reell sein soll. Die Bezeichnung »Zug komplexer Werte«, die \textsc{Gauss} hier anwendet, ist im Abschnitt [VI.] art. 2., oben S. 408, erklärt. Wir kommen so auf den Zusammenhang, der zwischen den Abschnitten [V.], [VI.] und [VII.] besteht, und damit auch zu den hier in Betracht kommenden geschichtlichen Feststellungen.
\par\medskip\noindent\rule{.24\linewidth}{.4pt}\par\smallskip
\noindent\hspace*{\parindent}\textsuperscript{*}) Es ist: \(\log_{10}\cos\varphi=9{,}8782827-10;\ \log_{10}(\frac12\cos\theta)=9{,}0536038-10\), also \(\theta=76^0\,55'\,20''\).
\par\endgroup
\ArithmeticHead{442}{L}{ANALYSIS. NACHLASS.}
